%% ma: L12-B14-0004
%% lop: 12
%% chuong: 5
%% bai: 14
%% dang: 12.14.4
%% loai: DS
%% muc: [NB, TH, VD, VD]
%% dap_an: "SĐSĐ"
%% nguon: "(NBV)-ĐỀ SỐ 4-MỖI NGÀY 1 ĐỀ THI 2025.pdf (D04, MỖI NGÀY 1 ĐỀ - Đề số 4), Phần II Câu 3"
%% kien_thuc: "Bài 7-8 (tọa độ vectơ, độ dài), Bài 14 (tích có hướng, vectơ pháp tuyến, góc giữa hai mặt phẳng)"
%% trang_thai: chinh-thuc
%% ghi_chu: "Đề gốc có ảnh minh họa (không cần). Đáp án gốc S Đ S Đ đúng (sympy: góc 11,31 độ; diện tích 688,37)."
\begin{ex}
Một kĩ sư xây dựng muốn thăm dò một mảnh đất dốc hình tứ giác $ABCD$ để sử dụng cho dự án phát triển bất động sản nhà ở. Kĩ sư sử dụng một thiết bị kinh vĩ tuyến điện tử để xác định vị trí bốn góc của khu đất. Toạ độ của bốn góc lần lượt là $A(15;25;9)$, $B(20;5;5)$, $C(-10;-10;2)$ và $D(-15;10;6)$ trong một hệ toạ độ $Oxyz$ mà gốc $O$ trùng với vị trí đứng của kĩ sư, đơn vị trên các trục lấy theo mét.
\choiceTF
{$\overrightarrow{AB}=(-5;20;4)$.}
{\True Mảnh đất có dạng hình bình hành.}
{Nếu mặt đất nằm ngang trùng với mặt phẳng toạ độ $Oxy$ thì độ nghiêng của mảnh đất so với mặt đất nằm ngang nhỏ hơn $10^\circ$.}
{\True Mảnh đất có diện tích khi làm tròn đến hàng đơn vị là $688\ \mathrm m^2$.}
\loigiai{a) $\overrightarrow{AB}=(5;-20;-4)$. Sai.
b) $\overrightarrow{DC}=(5;-20;-4)=\overrightarrow{AB}$ và $A,B,D$ không thẳng hàng nên $ABCD$ là hình bình hành. Đúng.
c) $\overrightarrow{AD}=(-30;-15;-3)$, $\left[\overrightarrow{AB},\overrightarrow{AD}\right]=(0;135;-675)=135(0;1;-5)$. Gọi $\varphi$ là góc giữa $(ABCD)$ và $(Oxy)$ (vectơ pháp tuyến $\vec k=(0;0;1)$): $\cos\varphi=\dfrac{5}{\sqrt{26}}\approx0{,}98$, suy ra $\varphi\approx11{,}3^\circ>10^\circ$. Sai.
d) $S=\left|\left[\overrightarrow{AB},\overrightarrow{AD}\right]\right|=135\sqrt{26}\approx688{,}4\approx688\ \mathrm m^2$. Đúng.}
\end{ex}
